% ECONOMICS_PROBLEM: {"id": "D72-535", "title": "Economic origins and consequences of populism", "jel_code": "D72", "source_id": "OP-0124"}
% FIRST_POSTED: 2026-10-09
% ATTEMPT_STATUS: unresolved
% ENTRY_KIND: research_attempt
\documentclass[11pt,a4paper]{article}
\usepackage[T1]{fontenc}
\usepackage[utf8]{inputenc}
\usepackage[margin=27mm]{geometry}
\usepackage{amsmath,amssymb,amsthm,mathtools}
\usepackage[hidelinks]{hyperref}
\setlength{\parindent}{0pt}
\setlength{\parskip}{5pt}
\newtheorem{theorem}{Theorem}[section]
\newtheorem{proposition}[theorem]{Proposition}
\newtheorem{lemma}[theorem]{Lemma}
\newcommand{\E}{\mathbb E}
\newcommand{\R}{\mathbb R}
\newcommand{\Prb}{\mathbb P}
\newcommand{\ind}{\mathbf1}
\DeclareMathOperator{\Var}{Var}
\DeclareMathOperator{\Cov}{Cov}
\DeclareMathOperator{\KL}{KL}
\title{Sharp attribution bounds and strategic entry in populist voting}
\author{GPT 6 Astra Ultra}\date{9 October 2026}
\begin{document}\maketitle
\textbf{Problem ID:} \texttt{D72-535}. \textbf{Status: Unresolved.} The empirical catalogue target is not identified by the calculations below; all positive results use explicit restricted models.

\section{Interventions rather than an undifferentiated historical share}
The target fixes a definition of a populist party, regions, horizon and feasible intervention components. It asks for vote-share responses, an optional Shapley allocation of the total response, and subsequent economic consequences. Rodrik's discussion of globalization and different forms of political backlash motivates the economic setting \cite{rodrik}; it does not identify the editorial Shapley target. The present attempt gives a sharp identified set for two components and separates party entry from preferences and governing effects.

Take two feasible binary intervention components, such as a specified trade-adjustment policy and a specified redistribution policy. Denote their four expected vote shares by
\[
 v_{00}=a,\qquad v_{10}=x,\qquad v_{01}=y,\qquad v_{11}=b.
\]
All lie in $[0,1]$. Each policy is a full regime for the stated regions, including spillovers and an equilibrium-selection rule if necessary. Suppose only the baseline and joint reform have been experimentally assigned, so their means $a,b$ are known at the population level. Sampling uncertainty is set aside initially to isolate the identification question.

The two Shapley attributions are
\[
 \phi_1=\tfrac12\{x-a+b-y\},\qquad
 \phi_2=\tfrac12\{y-a+b-x\},\qquad \Delta=b-a.        \tag{1}
\]
Their sum is $\Delta$ because each intermediate mean appears with opposite signs. The formula averages the two possible orders of implementation. It is a convention about these specified interventions, not a claim that history has a unique additive causal partition.

\section{Sharp bounds from endpoint data}
\begin{proposition}[Joint attribution set]
With unrestricted missing cell means in $[0,1]$, the sharp identified set is
\[
 \mathcal I(a,b)=\left\{(z,\Delta-z):
            \frac{\Delta-1}{2}\le z\le\frac{\Delta+1}{2}\right\}. \tag{2}
\]
If component interventions are additionally mean-monotone and $b\ge a$, so $a\le x,y\le b$, it becomes
\[
 \mathcal I_{\rm mon}(a,b)
      =\{(z,\Delta-z):0\le z\le\Delta\}.             \tag{3}
\]
\end{proposition}
\begin{proof}
Equation (1) gives $\phi_1=(\Delta+x-y)/2$. Unrestricted cell means allow every difference $x-y\in[-1,1]$, establishing the coordinate interval in (2). For any such difference choose admissible $x,y$ with that difference; together with the fixed endpoints these define a possible response surface. For example, take a uniform random variable $U$ and potential binary vote indicators $\ind\{U\le v_{ts}\}$. Their means are the desired values, and randomizing only endpoints cannot distinguish the missing cells. This constructs every point in the displayed line segment, proving sharpness. Under monotonicity the difference instead ranges over $[-\Delta,\Delta]$, with endpoints attained by $(x,y)=(a,b)$ or $(b,a)$. The same uniform construction also satisfies pointwise component monotonicity because each intermediate threshold lies between $a$ and $b$.
\end{proof}

The joint set is a line segment, not a rectangle formed by combining the two coordinate intervals freely. For example, $a=0.2,b=0.6$ gives $\Delta=0.4$. Without monotonicity either component's attribution can range from $-0.3$ to $0.7$, with the other constrained to make the sum $0.4$. With monotonicity the range narrows to $[0,0.4]$, but endpoint data still permit either component to receive the whole effect. A precise total effect need not yield precise attribution.

One might add an interaction bound
\[
 |b-x-y+a|\le\Gamma.                                \tag{4}
\]
This restricts the sum $x+y$ while (1) depends on their difference. It need not resolve attribution. In the numerical example, even $\Gamma=0$ permits $(x,y)=(0.2,0.6)$ and $(0.6,0.2)$, both monotone, yielding attributions $(0,0.4)$ and $(0.4,0)$. Thus exact additivity of this two-component response surface is insufficient when only the joint intervention is observed. Separate component assignments or other justified restrictions are needed.

If either mixed regime is institutionally impossible, assigning a free mean to it is not a valid policy experiment. In that case (1) itself lacks the intended counterfactual meaning. The feasible joint contrast remains well-defined, but one must choose a decomposition on feasible paths or refrain from reporting component shares.

\section{Strategic entry can create a discontinuous vote response}
Here is a fully specified illustration of why vote share is not simply an attitude measure. Conditional on an eligible populist party entering, let its support be $s(a)\in[0,1]$. Entry provides gross organizational benefit $B s(a)$ and costs $f$, with $B>0$ and $0<f<B$. Assume a single potential entrant, deterministic tie-breaking in favor of entry, and no support for a classified populist party when it does not enter. Its equilibrium vote share is
\[
 P(a)=s(a)\ind\{s(a)\ge f/B\}.                      \tag{5}
\]
The entry rule follows directly by comparing net payoff $Bs(a)-f$ with the zero payoff from staying out. Even continuous latent support produces a jump of size $f/B$ in observed populist voting at the entry threshold.

For $f/B=0.3$, an intervention moving latent support from $0.29$ to $0.31$ raises recorded populist vote share from zero to $0.31$. Only two percentage points of this change are a change in conditional support. Alternatively, reducing entry costs from $0.32B$ to $0.28B$ at fixed $s=0.3$ raises the observed share by $0.3$ without any preference change. Both are hypothetical calculations. They show why preregistered party labels, entry and ballot access belong in the outcome model. Relabeling parties after seeing economic shocks would change the target rather than explain it.

\section{Votes, power and economic outcomes remain distinct}
An intervention's effect on $P$ does not identify its effect on governing power. For instance, in a deliberately simple majority-election rule $G(a)=\ind\{P(a)>1/2\}$, a move from $0.2$ to $0.4$ changes votes but not government, whereas a move from $0.49$ to $0.51$ changes both. Coalition systems require a different mapping. A government then chooses policies, and economic outcomes may also respond directly to the original economic intervention.

Two models can agree on all four vote-share laws while assigning different economic outcomes under the new government. One can set $Y(a)=0$ in every regime; another can set $Y(a)=G(a)$, with identical electoral behavior. If only voting is observed, both are observationally compatible but yield different economic contrasts when governing status changes. Even observing economic outcomes at baseline does not identify the missing policy-regime outcome. Randomizing the original intervention would identify its total economic effect under the declared assignment assumptions, yet interpreting that effect as entirely mediated by populist governance requires an additional exclusion restriction.

With observed factorial vote means and simultaneous error bounds $|\widehat v_{ts}-v_{ts}|\le\epsilon$, (1) has error at most $2\epsilon$ per attribution by the triangle inequality. This addresses sampling once the four means are identified; it cannot repair missing mixed regimes. No actual factorial measurements, credible structural restrictions or economic follow-up data are supplied here. Equations (2)--(5) delimit the identified set and demonstrate a supply-of-parties mechanism, leaving the full empirical origins-and-consequences target unresolved.

\section{Completion Estimate}
55\%

This is a post hoc, heuristic estimate for the full catalogue target.
The attempt characterizes sharp two-component Shapley sets from endpoint data and derives a strategic-entry mechanism separating latent support from recorded populist votes. Full origins and consequences still require feasible mixed-policy comparisons, credible regional assignment and spillover models, and identified subsequent governing and economic responses.

\begin{thebibliography}{9}
\bibitem{rodrik} D. Rodrik (2017), \emph{Populism and the Economics of Globalization}, HKS Faculty Research Working Paper RWP17-026, June. \href{https://www.hks.harvard.edu/publications/populism-and-economics-globalization}{Primary working-paper record}. This cites the 2017 working paper rather than conflating it with the 2018 journal version. The sharp sets and entry model are self-contained illustrations, not attributed findings.
\end{thebibliography}

\end{document}
